\subsection{Mailman式共享求和：算术复用与散射成本}
\label{sec:8-mailman}
primary距离的符号码形式提示：共享位模式的候选可以复用求和。
Mailman提供有限字母矩阵--向量乘法的一般思路\cite{mailman}。
历史pilot固定Vogue、arXiv、OpenAI-3072候选流，在128 probe、batch 8/1,000
上测试GPU实现。基线复刻packed两表评分循环；替代方案包括原子桶累加、
私有直方图及4/6/8组宽，数学分数和输入编码不变。

六个单元中最快替代方案只有基线评分吞吐的0.187--0.581
（\cref{fig:8-negative}）。阶段profile把大多数插桩周期定位到桶累加。
list内偏置模式造成重复桶目标；私有直方图避免原子操作，却增加状态与共享内存
读改写。缩小组宽、绕过主要桶可以改善部分变体，但都未追上基线。
这支持停止这些实现与输入上的路线，不说明所有有限字母复用都无效：
胜出的两表基线已经沿另一方向复用坐标和。

pilot只写出primary分数，没有返回最终近邻。FP32/FP64检查揭示cutoff附近
数值差异，不等于完整流水线等价。以实测评分比$s$和估计可替换延迟占比$f$，
乐观替换模型为$1/[(1-f)+f/s]$；由于未集成，其预测减速只是模型。
辅助无损字典路线也在结构统计阶段停止，有限字节／工作节省没有支持GPU实现。
两条路线均未向主前沿贡献系统实测点。

\par\medskip\noindent\begin{minipage}{\linewidth}
\centering\captionsetup{hypcap=false}
\includegraphics[width=\linewidth]{ch8/negative-evidence.pdf}
\captionof{figure}{两种不同证据边界。左：RTX 5090、128 probe、Vogue、arXiv、OpenAI-3072和
两种batch下，最快Mailman评分器／packed表评分器速度，全部小于1。
右：arXiv的400查询FP64离线测试划分中，满足原JHQ recall目标的最小survivor数。
候选减少不是GPU提速；两图均由论文audit脚本从归档记录重新生成。}
\label{fig:8-negative}
\end{minipage}\par\medskip

\subsection{面向排序的编码：局部收益不等于通用替代}
\label{sec:8-rank}
排序导向研究在相同名义码长下，改变1-bit阈值或打分，尝试以更少survivor保留
有用近邻。主比较在导出候选上用FP64回放，划分为400校准、200验证、400最终测试。
方案包括似然权重、中位数阈值、学习排序损失、线性MMSE解码与list相关阈值。
较低的重建／内积MSE可能使极少数最近候选的顺序变差；改变primary电平还会
改变残差参照点，因此不同目标暴露不同失败模式。

arXiv测试集中，中位数阈值把达到原目标所需的最小被测survivor数从875降至375。
按同一严格标准，成功未推广到另外两个工作负载。共享旋转下的高斯直觉也不能
消除真实语料均值偏置，因为旋转保留均值向量的范数。
居中编码可以修复部分失配，但必须同时评价primary排序与重建分数上限。
因此停止的研究支持工作负载相关改进，而非普遍排序最优的新量化器。

后续仅median方案的GPU跟进是另一层证据：把改变后的阈值、电平及重新训练的
残差导入相同搜索流水线，分别扫预算，再按实测recall阈值比较。
在arXiv原$\alpha=100$目标，batch 8和1,000的比值约为1.40和1.21。
Vogue median前沿最高0.9630，未达到原工作点0.9644；OpenAI-3072基本持平。
这修正了“从未做GPU测试”的笼统表述，同时保留不采用通用替代的决定。
跟进使用独立索引状态、查询池和三轮协议，不能把其吞吐插入E1。
